\documentclass[UTF8,fontset=fandol,11pt]{ctexart}
\usepackage[a4paper,margin=24mm]{geometry}
\usepackage{amsmath,amssymb,amsthm,graphicx,tabularx,booktabs,hyperref}
\newtheorem{definition}{定义}[section]
\newtheorem{theorem}[definition]{定理}
\newtheorem{example}[definition]{例题}
\newtheorem{exercise}[definition]{练习}
\hypersetup{unicode=true,colorlinks=true,linkcolor=blue}
\title{Global repair}
\date{}
\begin{document}
\maketitle
\par\noindent 离线合成演示：不代表真实模型对本专题的教学验收。\par
\section{S01}\label{sec:S01}
% UNIT S01.U01
取足够大的正整数 N。 正整数无上界（阿基米德性质），因此这样的 N 存在。\par
% UNIT S01.U02
\begin{equation}1=1\label{eq:old-unchecked}\end{equation}\par
\end{document}
